\chapter{Kinetic Theory As A Controlled Limit}
\label{ch:thermal-kinetic-theory-controlled-limit}

\ConventionsMixed

Kinetic theory is an effective description of weakly disturbed many-particle
states in terms of phase-space densities.  Its QFT content lies in deriving
the quasiparticle variables, collision terms, and hydrodynamic moments from
real-time correlation functions under explicit scale assumptions.

\section{Quasiparticle Distribution}

Let \(a\) label a particle species with on-shell energy
\[
  E_a(\mathbf p)=\sqrt{\mathbf p^2+m_a^2}.
\]
A one-particle distribution function
\[
  f_a(t,\mathbf x,\mathbf p)
\]
is defined by a Wigner transform of a two-point function, projected onto a
quasiparticle pole and averaged over spacetime scales long compared with the
microscopic correlation length.  The projection is justified when the
spectral function has a narrow peak:
\[
  \Gamma_a(\mathbf p)\ll E_a(\mathbf p),
\]
where \(\Gamma_a\) is the damping width.

\section{Boltzmann Equation}

In the absence of external forces, the kinetic equation is
\[
  \left(\partial_t+\mathbf v_a(\mathbf p)\cdot\nabla_{\mathbf x}\right)
  f_a(t,\mathbf x,\mathbf p)
  =
  C_a[f](t,\mathbf x,\mathbf p),
\]
where
\[
  \mathbf v_a(\mathbf p)=\nabla_{\mathbf p}E_a(\mathbf p).
\]
For \(2\to2\) scattering, the collision term has the form
\[
\begin{aligned}
  C_a[f](p_1)
  &=
  \frac{1}{2E_1}
  \sum_{bcd}
  \int d\Pi_2\,d\Pi_3\,d\Pi_4\,
  (2\pi)^d\delta^{(d)}(p_1+p_2-p_3-p_4)
  |\mathcal M_{ab\to cd}|^2
\\
  &\quad\times
  \Big[
  f_c(p_3)f_d(p_4)(1\pm f_a(p_1))(1\pm f_b(p_2))
  -
  f_a(p_1)f_b(p_2)(1\pm f_c(p_3))(1\pm f_d(p_4))
  \Big],
\end{aligned}
\]
with invariant measure
\[
  d\Pi_i=\frac{d^{d-1}\mathbf p_i}{(2\pi)^{d-1}2E_i}.
\]
The sign is plus for bosons and minus for fermions.  The matrix element is
the on-shell scattering amplitude in the normalization of the scattering
chapters.

\section{Detailed Balance And Entropy}

The Bose-Einstein and Fermi-Dirac distributions
\[
  f_a^{\mathrm{eq}}(\mathbf p)
  =
  \frac{1}{\exp[\beta(E_a(\mathbf p)-\mu_a)]\mp1}
\]
annihilate the collision term when the chemical potentials obey the conserved
charge constraints for each reaction.  Define the kinetic entropy density
\[
  s_{\mathrm{kin}}
  =
  -\sum_a\int\frac{d^{d-1}\mathbf p}{(2\pi)^{d-1}}
  \left[
  f_a\log f_a
  \mp(1\pm f_a)\log(1\pm f_a)
  \right].
\]
Using the positivity of
\[
  (x-y)\log(x/y)\geq0
\]
for positive \(x,y\), the \(2\to2\) collision term gives nonnegative entropy
production.  Equality holds precisely when the gain and loss products agree
for every allowed collision.

\section{Hydrodynamic Moments}

The kinetic stress tensor and conserved current are
\[
  T^{\mu\nu}(x)
  =
  \sum_a\int d\Pi_p\,p^\mu p^\nu f_a(x,p),
  \qquad
  J^\mu_A(x)
  =
  \sum_a q_{A,a}\int d\Pi_p\,p^\mu f_a(x,p).
\]
Multiplying the Boltzmann equation by \(1\), \(p^\nu\), or a conserved charge
and integrating over momentum gives the hydrodynamic Ward identities when the
collision term conserves those quantities.  Constitutive relations are
obtained by solving the kinetic equation in a gradient expansion around local
equilibrium.

\section{Linearized Collision Operator}

Write
\[
  f_a=f_a^{\mathrm{eq}}+f_a^{\mathrm{eq}}(1\pm f_a^{\mathrm{eq}})\chi_a .
\]
To first order in \(\chi\), the collision term is a linear operator
\[
  C_a[f]=-\sum_b \mathcal L_{ab}\chi_b+O(\chi^2).
\]
The null space of \(\mathcal L\) is spanned by conserved quantities.  On the
orthogonal complement, positivity of entropy production makes
\(\mathcal L\) positive.  Transport coefficients are quadratic forms involving
\(\mathcal L^{-1}\) on this complement.

\section{Gauge Theories And Soft Scales}

In weakly coupled gauge theories, long-range gauge exchange produces soft
momentum transfers and collective plasma effects.  The kinetic description
must include screening, soft gauge fields, and collinear splitting processes
when they occur at the same order as \(2\to2\) scattering.  The separation of
hard quasiparticles, soft fields, and ultrasoft hydrodynamic modes is a
scale-by-scale effective theory construction, not a single Boltzmann equation
with arbitrary collision data.

\section{Controlled-Limit Ledger}

A kinetic derivation from QFT requires:
\[
\begin{array}{ll}
\text{spectral input} & \text{narrow quasiparticle peaks},\\
\text{scale separation} & \ell_{\mathrm{micro}}\ll \ell_{\mathrm{variation}},\\
\text{collision input} & \text{on-shell amplitudes with regulator and
medium corrections},\\
\text{conservation laws} & \text{collision invariants matching QFT Ward
identities},\\
\text{error estimate} & \text{bounds on neglected gradients and off-shell
terms}.
\end{array}
\]

\begin{openproblem}[QFT derivation of kinetic theory with errors]
\label{op:qft-kinetic-theory-error-derivation}
Derive relativistic kinetic theory from Schwinger-Keldysh QFT with explicit
error bounds, including quasiparticle projection, collision kernels,
screening, collinear processes in gauge theory, and the passage from kinetic
theory to hydrodynamic transport.
\end{openproblem}
